\chapter*{Abstract}
\addcontentsline{toc}{chapter}{Abstract}

Runtime monitoring can reveal failures and requirement violations that appear
only after a robot is deployed. ROS~2 exposes typed behaviour through topics,
services, and actions, but monitoring solutions differ in their observation
mechanisms, data formats, property checkers, and deployment assumptions.
Developers therefore repeat routine work for collection, conversion, routing,
transport, and evidence whenever they build a new monitoring setup.

This thesis presents a configurable monitoring infrastructure for ROS~2
applications. A domain analysis and feature model identify the commonalities and
variabilities of monitoring solutions. From this model, a reference architecture
separates ROS~2 observation, common data representation, transformation,
routing, checker adaptation, property checking, and result export. The
prototype collects topic messages, service request and response events, and
action feedback and status into a common JSON-compatible record. Converter and
verdict-service plug-ins isolate property-specific logic. A deployment
specification is converted into matching runtime YAML files for integrated or
split checking through file or MQTT transport.

Five experiments progress from a controlled topic property to unmodified Nav2,
two-robot correlation, and a deployment across Ubuntu, Raspberry Pi, and macOS.
In the reported runs, verdicts matched the transitions of the controlled
stimuli, the recorded inputs, the mission outcomes, or the simulated robot
positions. Identical local and remote verdict streams show that serialization
and MQTT transport did not alter the checking results in the evaluated
configurations. The generator and reuse analysis show that
observation, serialization, routing, transport, and evidence handling do not
have to be rewritten for each property. A runtime function mapping
also shows that the reference architecture can accommodate the principal
capabilities of four reviewed monitoring approaches through configuration,
plug-ins, and declared extension points.

The results support the proposed architecture and implementation within a
focused scope: passive application-level monitoring, selected ROS~2 interfaces,
up to two simulated robots, and one local network. The infrastructure provides a
reusable foundation on which additional observation mechanisms, checking
formalisms, transports, and operational safeguards can be added.
